\section{Perancangan Block Cipher}
Bab ini menjelaskan rancangan \algname{}: parameter, alur enkripsi dan dekripsi, pembentukan kunci
putaran, \textit{S-box}, permutasi dan rotasi, serta fungsi putaran berbasis sigma-delta.

\subsection{Parameter Umum}
\algname{} adalah \textit{block cipher} Feistel 16 putaran dengan blok dan kunci 128 bit. Fungsi
putarannya meminjam struktur modulator sigma-delta (Subbab~2.4).

\begin{table}[H]
  \centering
  \caption{Parameter \algname{}}
  \begin{tabular}{ll}
    \toprule
    Parameter & Nilai \\
    \midrule
    Ukuran blok & 128 bit (16 byte) \\
    Ukuran kunci utama & 128 bit (16 byte) \\
    Struktur & Jaringan Feistel dua cabang, masing-masing 64 bit \\
    Jumlah putaran & 16 \\
    Kunci putaran & 16 kunci, masing-masing 64 bit \\
    Ukuran \textit{word} pada fungsi $F$ & 32 bit \\
    \textit{S-box} & 8 bit, 256 entri, bijektif \\
    Permutasi bit & 64 bit, $i \mapsto 41\,i \bmod 64$ \\
    \bottomrule
  \end{tabular}
\end{table}

Ukuran blok 128 bit dipilih karena blok 64 bit rentan kolisi pada data besar (Subbab~2.1),
dan ukuran kunci 128 bit memenuhi syarat spesifikasi bahwa kunci utama minimal sepanjang
ukuran blok.

Tabel~\ref{tab:syarat} memetakan syarat spesifikasi ke komponen rancangan.

\begin{table}[H]
  \centering
  \caption{Pemenuhan syarat rancangan}\label{tab:syarat}
  \begin{tabular}{p{5.5cm}p{9cm}}
    \toprule
    Syarat & Komponen \\
    \midrule
    Substitusi (\textit{S-box}) & Tabel 256 entri pada fungsi $F$ dan \textit{key schedule} \\
    Transposisi (permutasi) & Permutasi bit 64 bit di akhir fungsi $F$ \\
    Operasi tambahan & Penjumlahan modulo $2^{32}$ dan rotasi \\
    Iterated cipher dengan kunci putaran & 16 putaran, kunci putaran dari \textit{key schedule} \\
    \textit{Confusion} & \textit{S-box} dan pencampuran kunci \\
    \textit{Diffusion} & Akumulator sigma-delta, rotasi, permutasi bit, dan struktur Feistel \\
    \bottomrule
  \end{tabular}
\end{table}

\subsection{Alur Enkripsi}
Blok plainteks 16 byte dibelah menjadi $L_0$ (8 byte kiri) dan $R_0$ (8 byte kanan), lalu
untuk $i = 1, \dots, 16$:
\begin{equation}
  L_i = R_{i-1}, \qquad R_i = L_{i-1} \oplus F(R_{i-1}, K_i).
\end{equation}
Setelah putaran ke-16, kedua separuh ditukar sekali lagi dan digabung menjadi cipherteks:
\begin{equation}
  C = R_{16} \,\|\, L_{16}.
\end{equation}

\begin{figure}[H]
  \centering
  \begin{tikzpicture}[
      box/.style={draw, minimum width=1.6cm, minimum height=0.7cm},
      xor/.style={draw, circle, inner sep=1pt},
      >=stealth]
    \node[box] (L0) at (0,0)    {$L_{i-1}$};
    \node[box] (R0) at (4,0)    {$R_{i-1}$};
    \node[box] (F)  at (4,-1.4) {$F$};
    \node[box] (K)  at (7,-1.4) {$K_i$};
    \node[xor] (X)  at (4,-2.8) {$\oplus$};
    \node[box] (L1) at (0,-4.6) {$L_i$};
    \node[box] (R1) at (4,-4.6) {$R_i$};
    \draw[->] (R0) -- (F);
    \draw[->] (K) -- (F);
    \draw[->] (F) -- (X);
    \draw[->] (L0.south) -- (0,-2.8) -- (X.west);
    \draw[->] (X) -- (R1);
    \draw[->] (R0.east) -- (5.2,0) -- (5.2,-3.7) -- (0,-3.7) -- (L1.north);
  \end{tikzpicture}
  \caption{Satu putaran Feistel pada \algname{}}\label{fig:round}
\end{figure}

\subsection{Alur Dekripsi}
Dekripsi memakai prosedur yang sama dengan enkripsi, dengan urutan kunci putaran dibalik
($K_{16}, K_{15}, \dots, K_1$). Pertukaran terakhir pada enkripsi membuat kedua prosedur
identik. Jika cipherteks $(R_{16}, L_{16})$ dimasukkan ke putaran pertama dekripsi dengan $K_{16}$:
\begin{equation}
  R_{16} \oplus F(L_{16}, K_{16}) = \bigl(L_{15} \oplus F(R_{15}, K_{16})\bigr) \oplus F(R_{15}, K_{16}) = L_{15},
\end{equation}
karena $L_{16} = R_{15}$. Keadaan setelah putaran ini adalah $(R_{15}, L_{15})$, yaitu keadaan
putaran ke-15 enkripsi dalam keadaan tertukar. Pola yang sama berlanjut hingga setelah 16 putaran
keadaannya $(R_0, L_0)$, dan pertukaran terakhir mengembalikan $(L_0, R_0)$, yaitu plainteks. Karena
yang dikombinasikan hanya hasil $F$ dan XOR, fungsi $F$ tidak perlu memiliki invers.

\subsection{Pembentukan \textit{Key Schedule}}
Kunci utama 128 bit dibelah menjadi dua \textit{word} 64 bit $a_0$ dan $b_0$. Untuk
$i = 1, \dots, 16$:
\begin{align*}
  a' &= S_{64}\bigl(\mathrm{rotl}_{64}(a_{i-1}, 13) + \varphi \cdot i \bmod 2^{64}\bigr), \\
  b' &= b_{i-1} \oplus \mathrm{rotl}_{64}(a', 29), \\
  (a_i, b_i) &= (b', a'), \\
  K_i &= a_i \oplus b_i,
\end{align*}
dengan $\varphi = \texttt{0x9E3779B97F4A7C15}$ (turunan \textit{golden ratio}), $\mathrm{rotl}_{64}$ rotasi
kiri 64 bit, dan $S_{64}$ substitusi tiap byte memakai \textit{S-box}.

Perancangan ini bertujuan agar:
\begin{itemize}
  \item tiap kunci putaran berbeda walau kunci utama berpola (mis.\ semua nol), melalui konstanta
    $\varphi \cdot i$ yang berubah tiap putaran;
  \item hubungan antar kunci putaran tidak linear terhadap kunci utama, melalui \textit{S-box} dan
    penjumlahan modulo $2^{64}$;
  \item setiap $K_i$ bergantung pada seluruh 128 bit kunci utama, karena $a$ dan $b$ saling
    mencampur lewat XOR dan pertukaran.
\end{itemize}
Nilai rotasi 13 dan 29 sengaja dipilih (bukan kelipatan 8), sehingga bit dalam satu byte terpecah ke
byte yang berbeda.

\subsection{\textit{S-Box}}
Tabel \textit{S-box} 8 bit dibangkitkan secara deterministik oleh program
\texttt{tools/gen\_sbox.cpp}: tabel identitas diacak dengan \textit{Fisher-Yates shuffle}
menggunakan PRNG xorshift64 dengan \textit{seed} 3141592653589793 (digit $\pi$), dan diulang
10.000 kali. Kandidat dengan titik tetap dibuang, kemudian dipilih kandidat dengan
\textit{nonlinearity} tertinggi dan, jika sama, \textit{differential uniformity} terendah.

\begin{table}[H]
  \centering
  \caption{Properti \textit{S-box} \algname{} dibanding pembanding}
  \begin{tabular}{lcc}
    \toprule
    Metrik & \algname{} & AES \\
    \midrule
    Bijektif & ya & ya \\
    Titik tetap $S(x) = x$ & 0 & 0 \\
    \textit{Nonlinearity} & 98 & 112 \\
    \textit{Differential uniformity} & 10 & 4 \\
    \bottomrule
  \end{tabular}
\end{table}

\textit{S-box} ini bersifat bijektif dan tanpa titik tetap, tetapi kualitas nonlinear dan
diferensialnya berada di kisaran \textit{S-box} acak (\textit{nonlinearity} 98 dan
\textit{differential uniformity} 10) dan di bawah \textit{S-box} AES yang dikonstruksi secara
aljabar. Ini keterbatasan pendekatan pencarian acak. Pada \algname{}, keamanan terhadap
kriptanalisis linear dan diferensial tidak bertumpu pada \textit{S-box} saja, tetapi pada
kombinasinya dengan akumulator, rotasi, permutasi, dan 16 putaran.

\subsection{Permutasi / Rotasi}\label{sec:perm}
\paragraph{Permutasi bit.} Bit pada posisi $i$ ($0 \le i < 64$) dipindah ke posisi
\begin{equation}
  \pi(i) = 41\,i \bmod 64.
\end{equation}
Karena $\gcd(41, 64) = 1$, $\pi$ bijektif. Pengali dipilih agar 8 bit dari satu byte tersebar ke
sebanyak mungkin byte tujuan, karena tugas utama permutasi di sini adalah memindahkan bit lintas
byte (\textit{S-box} bekerja per byte). Perbandingan beberapa pengali ganjil:

\begin{table}[H]
  \centering
  \caption{Penyebaran 8 bit satu byte ke byte tujuan untuk beberapa pengali}
  \begin{tabular}{ccc}
    \toprule
    Pengali & Byte tujuan berbeda (maks.~8) & Bit yang tidak berpindah \\
    \midrule
    37 & 7 & 4 \\
    \textbf{41} & \textbf{8} & 8 \\
    43 & 3 & 2 \\
    47 & 5 & 2 \\
    53 & 7 & 4 \\
    59 & 6 & 2 \\
    \bottomrule
  \end{tabular}
\end{table}

Pengali 41 adalah satu-satunya yang menyebarkan setiap byte ke 8 byte berbeda. Kelemahannya, 8 bit
(posisi kelipatan 8) tidak berpindah pada lapisan ini. Hal itu tidak fatal karena bit tersebut
tetap diolah \textit{S-box}, penjumlahan, dan rotasi pada putaran berikutnya.

\paragraph{Rotasi.} Rotasi kiri 32 bit dipakai di antara dua \textit{pass} sigma-delta
(Subbab~\ref{sec:fungsi}), dengan jumlah 7 untuk \textit{word} pertama dan 19 untuk \textit{word} kedua.
Rotasi diperlukan karena bit paling kanan hasil penjumlahan hanya bergantung pada bit paling
kanan operan (setara XOR), sehingga mudah ditebak. Rotasi memindahkannya ke posisi yang ikut
terkena \textit{carry} penjumlahan pada langkah berikutnya.

\subsection{Fungsi Putaran}\label{sec:fungsi}
Fungsi $F(R, K)$ menerima separuh blok $R$ (64 bit) dan kunci putaran $K$ (64 bit), dan
menghasilkan masker 64 bit. $R$ dan $K$ dipecah menjadi dua \textit{word} 32 bit,
$R = (w_0, w_1)$ dan $K = (k_0, k_1)$ dengan indeks 0 sebagai bagian atas.

\paragraph{Langkah sigma-delta.} Satu langkah memproses satu \textit{word} $x$ dengan kunci $k$,
membawa akumulator $a$ dan umpan balik $q_{\text{prev}}$:
\begin{align*}
  a &\leftarrow a + (x - q_{\text{prev}}) \bmod 2^{32}, \\
  q &= S_{32}(a \oplus k), \\
  q_{\text{prev}} &\leftarrow q,
\end{align*}
dengan $S_{32}$ substitusi tiap byte memakai \textit{S-box}. \textit{Output} langkah ini adalah $q$.
Nilai awal $a$ dan $q_{\text{prev}}$ adalah 0 pada tiap \textit{pass}.

\paragraph{Urutan fungsi $F$.}
\begin{enumerate}
  \item \textbf{\textit{Pass} maju}: $w_0 \leftarrow \mathrm{SD}(w_0, k_0)$, kemudian
    $w_1 \leftarrow \mathrm{SD}(w_1, k_1)$ dengan akumulator dan umpan balik dibawa dari $w_0$.
  \item \textbf{Rotasi}: $w_0 \leftarrow \mathrm{rotl}_{32}(w_0, 7)$,
    $w_1 \leftarrow \mathrm{rotl}_{32}(w_1, 19)$.
  \item \textbf{\textit{Pass} mundur}: $w_1 \leftarrow \mathrm{SD}(w_1, k_0)$, kemudian
    $w_0 \leftarrow \mathrm{SD}(w_0, k_1)$, kunci ditukar agar tiap \textit{word} terkena kedua
    kunci, dan akumulator dimulai dari 0 lagi.
  \item \textbf{Permutasi}: hasil $w_0 \,\|\, w_1$ dipermutasi dengan $\pi$ pada Subbab~\ref{sec:perm}.
\end{enumerate}

\paragraph{Hubungan dengan modulator sigma-delta.}
\begin{table}[H]
  \centering
  \caption{Pemetaan konsep sigma-delta ke fungsi $F$}
  \begin{tabular}{p{4.2cm}p{9.8cm}}
    \toprule
    Modulator sigma-delta & Fungsi $F$ \algname{} \\
    \midrule
    Akumulator & $a \leftarrow a + (x - q_{\text{prev}}) \bmod 2^{32}$ \\
    Kuantisasi & Diganti \textit{S-box} dengan XOR kunci (bijektif, tidak \textit{lossy}) \\
    Umpan balik & $q_{\text{prev}} \leftarrow q$ \\
    Aliran satu arah & Dijalankan maju lalu mundur agar dua arah \\
    (tidak ada) & Rotasi antar \textit{word} dan permutasi bit \\
    \bottomrule
  \end{tabular}
\end{table}

Rancangan ini hanya terinspirasi oleh modulator sigma-delta, bukan benar-benar \textit{copy one-to-one}. Kuantisasi 
aslinya membuang resolusi, yang cocok untuk DAC, tetapi tidak untuk cipher (menganggu \textit{integrity}). Sehingga, 
pada rancangan ini, diganti dengan \textit{S-box} bijektif. Yang dipertahankan adalah sifat bahwa tiap
\textit{output} bergantung pada seluruh \textit{input} sebelumnya melalui state yang dibawa, mirip mode CBC
atau CFB tempat satu blok bergantung pada blok sebelumnya.